\documentclass[11pt, a4paper]{article}

\usepackage{lipsum}

% Encoding and typography
\usepackage[utf8]{inputenc}
\usepackage[T1]{fontenc}
\usepackage{lmodern}
\usepackage{microtype}
\usepackage{parskip}

% Page layout
\usepackage[margin=1in]{geometry}
\usepackage[dvipsnames]{xcolor}

% Mathematics
\usepackage{amsmath}
\usepackage{amssymb}
\usepackage{amsthm}
\usepackage{mathtools}

% Algorithms
\usepackage{algorithm}
\usepackage{algpseudocode}
\usepackage[title]{appendix}

% Numerics section dependencies
\usepackage{graphicx}
\usepackage{adjustbox}
\usepackage{booktabs}
\usepackage{multicol}
\usepackage{caption}

% Theorem environments
\newtheoremstyle{theoremstyle}
  {\topsep}{\topsep} 
  {\itshape}{}       % Italic body
  {\bfseries}{.}     
  {5pt} 
  {\thmname{#1}\thmnumber{ #2}\thmnote{ (#3)}}

\newtheoremstyle{definitionstyle}
  {\topsep}{\topsep} 
  {\normalfont}{}    % Roman body
  {\bfseries}{.}     
  {5pt} 
  {\thmname{#1}\thmnumber{ #2}\thmnote{ (#3)}}

\theoremstyle{theoremstyle}
\newtheorem{theorem}{Theorem}[section]
\newtheorem{lemma}[theorem]{Lemma}
\newtheorem{proposition}[theorem]{Proposition}
\newtheorem{corollary}[theorem]{Corollary}

\theoremstyle{definitionstyle}
\newtheorem{definition}[theorem]{Definition}
\newtheorem{example}[theorem]{Example}
\newtheorem{exercise}[theorem]{Exercise}
\newtheorem{assumption}[theorem]{Assumption}
\newtheorem{remark}[theorem]{Remark}
\newtheorem{note}[theorem]{Note}

% Macros
\newcommand{\dd}{\mathop{}\!d}
\NewDocumentCommand{\DD}{ O{} O{\alpha} O{t} }{%
    \prescript{}{#1}{\rm D}^{#2}_{#3}%
}
\NewDocumentCommand{\II}{ O{} O{\alpha} O{t} }{%
    \prescript{}{#1}{\rm I}^{#2}_{#3}%
}
\NewDocumentCommand{\JJ}{ O{} O{\alpha} O{t} }{%
    \prescript{}{#1}{\rm J}^{#2}_{#3}%
}
\DeclareMathOperator{\PP}{\mathrm P}
\DeclareMathOperator{\EE}{\mathrm E}
\DeclareMathOperator{\linspan}{span}
\newcommand{\odv}[1]{\frac{d}{d#1}}
\newcommand{\pdv}[1]{\frac{\partial}{\partial#1}}
\DeclarePairedDelimiter{\norm}{\lVert}{\rVert}
\DeclarePairedDelimiter{\abs}{\lvert}{\rvert}

% Hyperlinks and references
\usepackage{hyperref}
\hypersetup{
    colorlinks=true,
    linkcolor=blue!70!black,
    citecolor=green!70!black,
    urlcolor=blue!70!black,
}
\usepackage[style=numeric]{biblatex}
\addbibresource{references.bib}

\title{An Operational Matrix Approach to M\"untz-Legendre Deep Neural Networks as a Solution to Nonlinear Fractional Stochastic Differential Equations}
\date{}
\author{}

\begin{document}

\maketitle

\section{Introduction}

\subsection{Notation and Preliminaries}


\begin{definition}[Fractional operators]\label{def:fp}
    Let $\alpha \in (1/2, 1)$ denote the fractional order of the following operators:
    \begin{enumerate}
        \item \textbf{Riemann--Liouville fractional integral:}
        \begin{equation} \II f(t) = \frac{1}{\Gamma(\alpha)} \int_0^t (t-\tau)^{\alpha-1} f(\tau) \dd{\tau}. \end{equation}
        \item \textbf{Caputo fractional derivative:}
        \begin{equation} \DD f(t) = \frac{1}{\Gamma(1-\alpha)} \int_0^t (t-\tau)^{-\alpha} f'(\tau) \dd{\tau}. \end{equation}
        \item \textbf{Stochastic Riemann--Liouville fractional integral:}
        \begin{equation} \JJ f(t) = \frac{1}{\Gamma(\alpha)} \int_0^t (t-\tau)^{\alpha-1} f(\tau) \dd{B_\tau}. \end{equation}
    \end{enumerate}
\end{definition}

\begin{definition}[Inner product]\label{def:ip}
    All orthogonality relations and projections in this work are defined with respect to the standard (unweighted) \( L^2 \) inner product on \( [0,1] \) and associated norm:
    \begin{equation} \langle f, g \rangle = \int_0^1 f(t) g(t) \dd{t}, \quad \norm{f} = \sqrt{\langle f, f \rangle}. \end{equation}
\end{definition}

\begin{definition}[Mittag-Leffler function]\label{def:ml}
    For $\alpha, \beta > 0$, the two-parameter Mittag-Leffler function is defined by the series
    \begin{equation} E_{\alpha,\beta}(z) = \sum_{k=0}^{\infty} \frac{z^k}{\Gamma(\alpha k + \beta)}, \quad z \in \mathbb C, \end{equation}
    with $E_\alpha(z) \coloneqq E_{\alpha,1}(z)$.
\end{definition}

% In typical experimental setups, we utilize $\hat{m} = 8$, $N_q = 64$, $n^* = 200$, and $r = 3$.

\subsection{Model Problem}\label{sec:model_prob}

Let \( (\Omega, \mathcal F, \{\mathcal F_t\}_{t \in [0,1]}, \PP) \) be a complete filtered probability space satisfying the usual conditions, where the filtration is generated by a standard one-dimensional Brownian motion on \([0, 1]\) and augmented by all \( \PP \)-null sets.

Consider the Caputo fractional stochastic differential equation (CFSDE) of order \( \alpha \) for a stochastic process \( y : \Omega \times [0, 1] \to \mathbb{R} \):
\begin{equation} \DD y(t) = b(t, y(t)) + \sigma(t, y(t))\dd{B_t}, \quad t \in [0, 1], \end{equation}
subject to the initial condition \( y(0) = y_0 \), where \( y_0 \) is an \( \mathcal F_0 \)-measurable random variable satisfying \( \EE [y_0^2] < \infty \).

\begin{assumption}[Well-posedness]\label{as:wp}
    The drift coefficient $b\colon [0,1] \times \mathbb{R} \to \mathbb{R}$ and diffusion coefficient $\sigma\colon [0,1] \times \mathbb{R} \to \mathbb{R}$ are Borel measurable functions that satisfy the following conditions:
    \begin{enumerate}
        \item \textbf{Lipschitz continuity.} There exists a constant $L > 0$ such that for all $t \in [0,1]$ and $x, y \in \mathbb{R}$,
        \begin{equation} |b(t,x) - b(t,y)| + |\sigma(t,x) - \sigma(t,y)| \leq L|x - y|. \end{equation}
        \item \textbf{Linear growth bound.} There exists a constant $K > 0$ such that for all $t \in [0,1]$ and $x \in \mathbb{R}$,
        \begin{equation} |b(t,x)|^2 + |\sigma(t,x)|^2 \leq K^2(1 + |x|^2). \end{equation}
    \end{enumerate}
    Under these conditions, the CSFDE admits a unique strong solution $y$, which is an $\{\mathcal{F}_t\}$-adapted stochastic process such that $y \in L^2(\Omega; C([0,1]; \mathbb{R}))$.
\end{assumption}

\subsection{Network Architecture}\label{sec:net_arch}
Define a M\"untz sequence $\Lambda = \{\lambda_0, \lambda_1, \lambda_2, \dots, \lambda_{\hat m}\}$, where $\lambda_0 = 0$ and the subsequent terms $\lambda_k$ are chosen as multiples of the fractional order $\alpha$. The M\"untz-Legendre polynomials
\begin{equation} \mathbf M^\Lambda (t) = [M_1^\Lambda(t), M_2^\Lambda(t), \dots, M_{\hat m}^\Lambda(t)]^\top \end{equation}
are orthogonalized over the interval $[0, 1]$.

The neural network $N_\theta(t)$ acts as a universal approximator for a specific sample path $y(t, \omega)$. We define this architecture by $L$ layers ($\mathcal D_i$ for $i \in \{0, 1, \dots, L\}$), incorporating an explicit fractional feature map:
\begin{subequations}
\begin{align}
    \mathcal D_0 &= t, && \text{(Input layer)} \\
    \mathcal D_1 &= \mathbf M^\Lambda(\mathcal D_0), && \text{(Fractional feature layer)} \\
    \mathcal D_i &= f(\mathbf W^{(i)}\mathcal D_{i-1} + \mathbf b^{(i)}), \quad i \in \{2, \dots, L-1\}, && \text{(Hidden layers)} \\
    N_\theta (t) &= \mathbf W^{(L)}\mathcal D_{L-1} + \mathbf b^{(L)} && \text{(Output layer)},
\end{align}
\end{subequations}
where $f$ is a nonlinear activation function, and $\theta = \{\mathbf W^{(i)}, \mathbf b^{(i)}\}_{i=1}^L$ are the trainable parameters of the network.

\section{Volterra-Integral Reformulation via Operational Matrices}\label{sec:volterra}
To solve the model problem, we propose a unified methodology based on operational matrices that bypasses the need for dynamic numerical integration or the approximation of inherently nondifferentiable white noise. 

Applying the fractional Riemann-Liouville integral operator on both sides of the CSFDE, we obtain its equivalent integrated form, avoiding direct evaluation of the Caputo operator:
\begin{equation} y(t) = y_0 + \frac{1}{\Gamma(\alpha)} \int_0^t (t-\tau)^{\alpha-1} b(\tau, y(\tau)) \dd{\tau} + \frac{1}{\Gamma(\alpha)} \int_0^t (t-\tau)^{\alpha-1} \sigma(\tau, y(\tau)) \dd{B_\tau}. \end{equation}

To eliminate numerical integration during the training process, the nonlinear drift function $b(t, N_\theta(t))$ and the diffusion function $\sigma(t, N_\theta(t))$ are mapped onto the M\"untz-Legendre space using learnable projection vectors $\theta_b$ and $\theta_\sigma$:
\begin{equation}
\begin{aligned}
    b(t, N_\theta(t)) &\approx \theta_b^\top \mathbf{M}^\Lambda(t), \\
    \sigma(t, N_\theta(t)) &\approx \theta_\sigma^\top \mathbf{M}^\Lambda(t).
\end{aligned}
\end{equation}

\subsection{Stochastic and Deterministic Operational Matrices}\label{sec:stoch_det_op}
By substituting these projections back into the integral equation, the continuous integral transformations are mapped algebraically through fixed operational matrices:
\begin{equation}
\begin{aligned}
    \frac{1}{\Gamma(\alpha)} \int_0^t (t-\tau)^{\alpha-1} b(\tau, y(\tau)) \dd{\tau} &\approx \theta_b^\top \mathbb P_\alpha(t) \mathbf{M}^\Lambda(t), \\
    \frac{1}{\Gamma(\alpha)} \int_0^t (t-\tau)^{\alpha-1} \sigma(\tau, y(\tau)) \dd{B_\tau} &\approx \theta_\sigma^\top \mathbb S_\alpha \mathbf{M}^\Lambda(t),
\end{aligned}
\end{equation}
where $\mathbb P_\alpha(t)$ is the deterministic fractional integration operational matrix, and $\mathbb{S}_\alpha$ is the weakly singular stochastic operational matrix evaluating the fractional stochastic integral.

\begin{remark}[Restriction on $\boldsymbol \alpha$]\label{rem:alpha}
    The construction of $\mathbb{S}_\alpha$ requires the kernel $(t-\tau)^{\alpha-1}$ to be square-integrable with respect to $\tau$ over $[0,t]$. Since
    \begin{equation} \int_0^t (t-\tau)^{2(\alpha-1)} \dd{\tau} = \frac{t^{2\alpha-1}}{2\alpha - 1}, \end{equation}
    finiteness demands $2\alpha - 1 > 0$, i.e., $\alpha > 1/2$. This condition ensures that the fractional It\^o integral
    \begin{equation} \int_0^t (t-\tau)^{\alpha-1} g(\tau) \dd{B_\tau} \end{equation}
    is well-defined as an element of $L^2(\Omega)$ for adapted, square-integrable integrands $g$.
%    Extending the methodology to $\alpha \in (0, 1/2]$ would require alternative integral formulations (e.g., a Skorokhod-type framework or regularization of the singular kernel) and is beyond the scope of the present work.
\end{remark}

\subsubsection{Constructing the Deterministic Operational Matrix}\label{sec:det_op}
Following the analytical framework established in \cite[Section 3.3]{SINGH2016195}, the operational matrix of integration is constructed by mapping the weakly singular integration kernel to the underlying basis space.

Apply the Riemann--Liouville fractional integral operator to a single M\"untz monomial component $t^{\lambda_j}$ (where $\lambda_j \in \Lambda$ represents a real-valued exponent). We evaluate the exact analytical scaling via the beta function formulation:
\begin{equation} \II (t^{\lambda_j}) = \frac{1}{\Gamma(\alpha)} \int_0^t (t-\tau)^{\alpha-1} \tau^{\lambda_j} \dd{\tau} = \frac{\Gamma(\lambda_j + 1)}{\Gamma(\lambda_j + \alpha + 1)} t^{\lambda_j + \alpha}. \end{equation}

Extending this element-wise integration to the underlying M\"untz monomial vector $\mathbf{X}(t) = [t^{\lambda_0}, t^{\lambda_1}, \dots, t^{\lambda_{\hat m}}]^\top$, we can algebraically isolate the time-dependent power-shift factor $t^\alpha$:
\begin{equation} \II \mathbf{X}(t) = t^\alpha \mathbf{D}^\Lambda_\alpha \mathbf{X}(t), \end{equation}
where $\mathbf{D}^\Lambda_\alpha$ is a static, purely deterministic diagonal matrix containing the ratio of Gamma functions:
\begin{equation} \mathbf{D}^\Lambda_\alpha = \mathrm{diag}\left( \frac{\Gamma(\lambda_0 + 1)}{\Gamma(\lambda_0 + \alpha + 1)}, \frac{\Gamma(\lambda_1 + 1)}{\Gamma(\lambda_1 + \alpha + 1)}, \dots, \frac{\Gamma(\lambda_{\hat m} + 1)}{\Gamma(\lambda_{\hat m} + \alpha + 1)} \right). \end{equation}

Because the orthogonal M\"untz-Legendre polynomial basis elements $\mathbf{M}^\Lambda(t)$ are constructed as linear combinations of these raw monomials, they are bound via a constant, invertible transformation matrix $\mathbf{C}$:
\begin{equation} \mathbf{M}^\Lambda(t) = \mathbf{C} \mathbf{X}(t) \Rightarrow \mathbf{X}(t) = \mathbf{C}^{-1} \mathbf{M}^\Lambda(t). \end{equation}

Utilizing the linearity of the fractional integration operator, we evaluate the transformation on the network's structural feature layer $\mathbf{M}^\Lambda(t)$ by pulling out the constant coefficient matrix $\mathbf{C}$ and substituting the monomial mappings:
\begin{equation}
\begin{aligned}
    \II \mathbf{M}^\Lambda(t) &= \II \left( \mathbf{C} \mathbf{X}(t) \right) \\
    &= \mathbf{C} \left( \II \mathbf{X}(t) \right) \\
    &= \mathbf{C} \left( t^\alpha \mathbf{D}_\alpha \mathbf{X}(t) \right) \\
    &= t^\alpha \mathbf{C} \mathbf{D}_\alpha \left( \mathbf{C}^{-1} \mathbf{M}^\Lambda(t) \right) \\
    &= \left( t^\alpha \mathbf{C} \mathbf{D}_\alpha \mathbf{C}^{-1} \right) \mathbf{M}^\Lambda(t).
\end{aligned}
\end{equation}

By matching this result with the definition of the operational matrix of integration, $\II \mathbf{M}^\Lambda(t) \approx \mathbb P_\alpha(t) \mathbf{M}^\Lambda(t)$, the operational matrix (exact up to truncation at degree $\hat{m}$) is explicitly isolated as:
\begin{equation} \mathbb P_\alpha(t) = t^\alpha \mathbf{C} \mathbf{D}_\alpha \mathbf{C}^{-1}. \end{equation}

\subsubsection{Constructing the Stochastic Operational Matrix}\label{sec:stoch_op}
Adapting the spectral projection framework established in \cite[Section 3]{gupta2025operational} for stochastic Volterra integral equations with weakly singular kernels, we construct the stochastic operational matrix \( \mathbb{S}_\alpha \). Let \( \mathbf{I}(t) = [I_1(t), I_2(t), \dots, I_{\hat m}(t)]^\top \) represent the time-dependent stochastic integral vector obtained by applying the weakly singular fractional Itô integral operator to each element of the M\"untz-Legendre basis vector \( \mathbf{M}^\Lambda(t) \):
\begin{equation} I_i(t) = \frac{1}{\Gamma(\alpha)} \int_0^t (t-\tau)^{\alpha-1} M_i^\Lambda(\tau) \dd{B_\tau}, \quad i \in \{1, 2, \dots, \hat m\}. \end{equation}
For a chosen random sample path realization \( (\omega) \), the vector \( \mathbf{I}(t) \) represents a continuous, non-smooth, stochastic process over the temporal domain \( t \in [0,1] \). To eliminate the need for repeating this stochastic convolution at every training epoch, we approximate \( \mathbf{I}(t) \) by projecting it back onto the underlying orthogonal M\"untz-Legendre space via an operational transformation:
\begin{equation} \mathbf{I}(t) \approx \mathbb{S}_\alpha \mathbf{M}^\Lambda(t), \end{equation}
where \( \mathbb{S}_\alpha \in \mathbb{R}^{\hat m \times \hat m} \) is the static stochastic operational matrix of integration for the given sample path. Extracting the \( i \)-th row of this matrix system yields the expansion for a single integral component:
\begin{equation} I_i(t) \approx \sum_{j=1}^{\hat m} (\mathbb{S}_\alpha)_{ij} M_j^\Lambda(t). \end{equation}
To isolate the constant matrix coefficients \( (\mathbb{S}_\alpha)_{ij} \), we take the standard \( L^2 \) inner product of both sides with respect to an arbitrary basis polynomial \( M_k^\Lambda(t) \) over the domain \( [0, 1] \):
\begin{equation} \langle I_i(t), M_k^\Lambda(t) \rangle = \sum_{j=1}^{\hat m} (\mathbb{S}_\alpha)_{ij} \langle M_j^\Lambda(t), M_k^\Lambda(t) \rangle. \end{equation}
Exploiting the strict orthogonality of the M\"untz-Legendre polynomials, the inner product \( \langle M_j^\Lambda(t), M_k^\Lambda(t) \rangle \) vanishes for all \( j \neq k \). This reduces the summation to a single non-zero term, allowing us to solve explicitly for each entry of the matrix:
\begin{equation}\label{eq:25} (\mathbb{S}_\alpha)_{ik} = \frac{\langle I_i(t), M_k^\Lambda(t) \rangle}{\norm{M_k^\Lambda}^2} = \frac{1}{\norm{M_k^\Lambda}^2} \int_0^1 I_i(t) M_k^\Lambda(t) \dd{t}. \end{equation}
By writing this relationship in full matrix form, the stochastic operational matrix can be expressed compactly as:
\begin{equation} \mathbb{S}_\alpha = \mathbf{\Xi} \boldsymbol{\mu}^{-1}, \end{equation}
where \( \boldsymbol{\mu} = \text{diag}(\norm{M_1^\Lambda}^2, \norm{M_2^\Lambda}^2, \dots, \norm{M_{\hat m}^\Lambda}^2) \) is the deterministic diagonal matrix of the squared norms of the basis elements, and \( \mathbf{\Xi} \in \mathbb{R}^{\hat m \times \hat m} \) is the path-dependent stochastic cross-covariance matrix whose elements are defined by the definitive double integral:
\begin{equation} \Xi_{ik} = \frac{1}{\Gamma(\alpha)} \int_0^1 \left[ \int_0^t (t-\tau)^{\alpha-1} M_i^\Lambda(\tau) \dd{B_\tau} \right] M_k^\Lambda(t) \dd{t}. \end{equation}
Because \( \mathbf{\Xi} \) depends solely on the pre-sampled Brownian motion path \( (\omega) \) and the fixed basis profiles, it is computed exactly once via fine-mesh initialization before the neural network training loop commences, remaining an absolute constant tensor during optimization.

\subsection{Loss Function Formulation}\label{sec:loss_form}
Rather than utilizing a graded mesh, discretization is performed over a collocation set $\{t_i\}_{i=1}^{N_q}$ using shifted Chebyshev roots
\begin{equation}
    \frac{1}{2}\left (1 - \cos \frac{(2i-1)\pi}{2N_q}\right ), \quad i \in \{1, 2, \dots , N_q\}, 
\end{equation}
to optimize global polynomial interpolation and maintain a well-conditioned loss landscape.

The total optimization objective minimizes the integrated SDE structural residual simultaneously with the learned projection consistency requirements:
\begin{equation} \mathcal L_{\text{total}}(\theta_{\text{SDE}}, \theta_b, \theta_\sigma) = \mathcal L_{\text{SDE}} + \lambda_b \mathcal L_b + \lambda_\sigma \mathcal L_\sigma, \end{equation}
where $\lambda_b$ and $\lambda_\sigma$ are balancing regularizers, and the loss components are defined explicitly as:
\begin{subequations}
\begin{align}
    \mathcal{L}_{\text{SDE}} &= \frac{1}{N_q} \sum_{i=1}^{N_q} \norm{ N_\theta(t_i) - y_0 - \theta_b^\top \mathbb P_\alpha(t_i) \mathbf{M}^\Lambda(t_i) - \theta_\sigma^\top \mathbb S_\alpha \mathbf{M}^\Lambda(t_i) }^2, \label{eq:29a} \\
    \mathcal{L}_{b} &= \frac{1}{N_q} \sum_{i=1}^{N_q} \norm{ \theta_b^\top \mathbf{M}^\Lambda(t_i) - b(t_i, N_\theta(t_i)) }^2, \\
    \mathcal{L}_{\sigma} &= \frac{1}{N_q} \sum_{i=1}^{N_q} \norm{ \theta_\sigma^\top \mathbf{M}^\Lambda(t_i) - \sigma(t_i, N_\theta(t_i)) }^2.
\end{align}
\end{subequations}

\section{Convergence Analysis}
We present the theoretical convergence proof for the proposed framework, demonstrating that the neural network approximation converges to the true solution of the CSFDE under appropriate conditions. Throughout this section, we fix a realized Brownian motion sample path \( B(t) \).

\subsection{Preliminaries}
We first establish preliminary results on the error bounds for the various operational matrices utilized in the proposed method.

\begin{lemma}\label{lem:3.1}
    For every \( t \in [0, 1] \) and every \( \theta_b \in \mathbb R^{\hat m} \), the identity
    \begin{equation}
        \theta^\top_b \mathbb P_{\alpha}(t)\mathbf M^\Lambda(t) = \II \left [\theta^\top_b \mathbf M^\Lambda \right ](t).
    \end{equation}
    holds exactly.
\end{lemma}

The proof for this lemma is essentially a repeat of the construction of the matrix in \S\ref{sec:det_op}. As a consequence of the lemma, every source of deterministic error in \( \mathcal L_{\text{SDE}} \) is only due to the error in approximation of \( \theta_b^\top \mathbf M^\Lambda(t) \) as the true drift \( b(t, N_\theta(t)) \) (that is, through \( \mathcal L_b \)) and not the operational matrix itself.

\begin{lemma}\label{lem:3.2}
    Let \( \mathbf I(t) \) be as defined in \S\ref{sec:stoch_op}. For any \( \theta_\sigma \in \mathbb R^{\hat m} \) such that \( \norm{\theta_\sigma} \leq \Theta \), it holds that
    \begin{equation}
        \EE \left [ (\theta_\sigma^\top \mathbf I(t) - \theta_\sigma^\top \mathbb S_\alpha \mathbf M^\Lambda(t))^2 \right ] \leq \Theta^2 \delta_S(\hat m)^2
    \end{equation}
    where \( \delta_S(\hat m) = O(\hat m^{-(\alpha - 1/2) + \epsilon}) \) for every \( \epsilon > 0 \).
\end{lemma}
\begin{proof}
    As constructed earlier in \S\ref{sec:stoch_op}, \( \mathbb S_\alpha \mathbf M^\Lambda (t) \) is the \( \hat m \) degree \( L^2[0, 1] \) orthogonal projection of the stochastic integral vector \( \mathbf I(t) \) onto \(\mathrm{span} \{M_1^\Lambda, \dots , M_{\hat m}^\Lambda\}\) with coefficients described by Equation~\ref{eq:25}.

    To bound the error of this polynomial projection, we must establish the regularity of a fixed component \( I_i \). In accordance with Remark~\ref{rem:alpha}, since \( \alpha > 1/2 \), the fractional kernel is square-integrable. Consequently, by the Kolmogorov continuity theorem for fractional Volterra processes, the sample paths of the stochastic integral \( I_i \) exhibit H\"older continuity of order \( \beta < \alpha - 1/2 \). 

    Using standard polynomial approximation theory \cite[\S 7.6]{devore1993constructive}, the \( L^2[0,1] \) orthogonal projection error for a H\"older-\(\beta\) function onto Müntz polynomials of degree \( \hat{m} \) decays as \( O(\hat{m}^{-\beta}) \). Substituting the sharp regularity exponent \( \beta = \alpha - 1/2 - \epsilon \), the error for each individual component is bounded by:
    \begin{equation}
        \norm{I_i - (\mathbb{S}_\alpha \mathbf{M}^\Lambda)_i} = O(\hat{m}^{-(\alpha - 1/2) + \epsilon})
    \end{equation}
    
    To bound the total projection error for the term \( \theta_\sigma^\top \mathbf{I}(t) \), we apply the Cauchy-Schwarz inequality across the \( \hat{m} \) components. Given the parameter bound \( \norm{\theta_\sigma} \leq \Theta \), taking the expectation yields:
    \begin{equation}
        \EE \left[ \left ( \theta_\sigma^\top \mathbf{I}(t) - \theta_\sigma^\top \mathbb S_\alpha \mathbf M^\Lambda(t) \right )^2 \right] \le \Theta^2 \delta_S(\hat{m})^2
    \end{equation}
    where \( \delta_S(\hat{m}) = O(\hat{m}^{-(\alpha - 1/2) + \epsilon}) \).
\end{proof}

We now state an assumption on the training outputs.

\begin{assumption}\label{as:3.3}
    There exist real constants \(\Theta\) and \(\Theta'\) independent of \(\hat m\) and \(N_q\) such that every trained network in the family under consideration satisfies \(\norm{\theta_b}, \norm{\theta_\sigma} \leq \Theta\), and the mappings \(t \mapsto N_\theta(t)\), \(t \mapsto \theta_b^\top \mathbf M^\Lambda(t)\), and \(t \mapsto \theta_\sigma^\top \mathbf M^\Lambda (t)\) are H\"older of order \(p \in (0, 1]\) with H\"older seminorm bounded by \(\Theta'\).
\end{assumption}
This assumption holds for the architecture in \S\ref{sec:net_arch}. \(N_\theta\) has a fixed depth with smooth activation, and finite trained weights on a compact domain. Furthermore, \(\mathbf M^\Lambda\) is smooth on \((0, 1]\). We state this assumption explicitly since it enables moving from control at \(N_q\) discrete points to control on all of \([0, 1]\).

\begin{lemma}\label{lem:3.4}
    Under Assumption~\ref{as:3.3}, let \( g \) be a H\"older function of order \( p \) satisfying
    \begin{equation}
        \frac 1{N_q} \sum_{i=1}^{N_q} g(t_i)^2 \leq \eta^2
    \end{equation}
    at the \( N_q \) shifted Chebyshev collocation points of \S\ref{sec:loss_form}. Then there exists a constant \(C_p\) such that
    \begin{equation}
        \sup_{t \in [0, 1]} \abs{g(t)} \leq \eta + C_p\Theta'N_q^{-p}.
    \end{equation}
\end{lemma}
\begin{proof}
Let \( P_{N_q}(t) \) denote the unique polynomial of degree at most \( N_q - 1 \) that interpolates \( g(t) \) at the shifted Chebyshev nodes \( \{t_i\}_{i=1}^{N_q} \). By definition, \( P_{N_q}(t_i) = g(t_i) \) for all \( 1 \le i \le N_q \).

Since \( g \) is H\"older of order \( p \) with a H\"older seminorm strictly bounded by \( \Theta' \), Jackson's theorems \cite[\S2.9, \S7.6]{devore1993constructive} guarantees the existence of a best-approximation polynomial \( P^*(t) \) of degree \( N_q - 1 \) whose uniform error satisfies.
\begin{equation}
    \norm{g - P^*}_\infty \le C_1 \Theta' N_q^{-p}
\end{equation}
where \( C_1 \) is a constant dependent only on \( p \).

The interpolant \( P_{N_q}(t) \) is related to the best approximation \( P^*(t) \) via the Lebesgue constant \( \Lambda_{N_q} \). By Lebesgue's lemma \cite[\S2.4]{devore1993constructive}, the uniform interpolation error is bounded as
\begin{equation}
    \norm{g - P_{N_q}}_\infty \le (1 + \Lambda_{N_q}) \norm{g - P^*}_\infty.
\end{equation}
For shifted Chebyshev nodes, the Lebesgue constant grows at a near-optimal, logarithmic rate \( \Lambda_{N_q} = O(\log N_q) \) \cite{Bos2012}, as the Lebesgue constant is strictly invariant under a linear domain transformation \cite{Ibrahimoglu2016}. This logarithmic factor is dominated by the polynomial decay \( N_q^{-p} \). Consequently, the maximum pointwise deviation of the interpolant from the true function across the continuous domain is bounded as
\begin{equation}
    \sup_{t \in [0,1]} |g(t) - P_{N_q}(t)| \le C_2 \Theta' N_q^{-p}.
\end{equation}

At the collocation nodes, the root-mean-square magnitude of the function is explicitly bounded by hypothesis:
\begin{equation}
    \left ( \frac{1}{N_q} \sum_{i=1}^{N_q} g(t_i)^2 \right )^{1/2} \le \eta.
\end{equation}
Because \( P_{N_q}(t) \) interpolates \( g(t) \) exactly at these nodes, its discrete magnitudes are controlled by \( \eta \). To establish the global bound, we apply the triangle inequality for any arbitrary continuous point \( t \in [0,1] \):
\begin{equation}
    \abs{g(t)} \le \abs{P_{N_q}(t)} + \abs{g(t) - P_{N_q}(t)}.
\end{equation}
Substituting the uniform bound for the continuous deviation and the discrete bound \( \eta \) controlling the interpolant yields the final inequality:
\begin{equation}
    \sup_{t \in [0,1]} \abs{g(t)} \le \eta + C_p \Theta' N_q^{-p}
\end{equation}
where \( C_p \) is a synthesized constant dependent only on \( p \).
\end{proof}

We now state one final result: the fractional Gronwall inequality \cite{YE20071075}.

\begin{lemma}\label{lem:3.5}
    Let \( \beta \in (0,1) \), \(K, \Phi \geq 0\), and \(h : [0, 1] \to [0, \infty)\) satisfy
    \begin{equation}
        h(t) \leq \Phi + K\int_0^t (t - \tau)^{\beta-1}h(\tau)\dd{\tau}, \quad t \in [0, 1].
    \end{equation}
    Then, it follows that \( h(t) \leq \Phi E_\beta(K\Gamma(\beta)t^{\beta}) \) for every \(t \in [0, 1]\), where \(E_\beta\) is the Mittag-Leffler function (Definition~\ref{def:ml}).
\end{lemma}

\subsection{Proof of Convergence}
We are now ready to tackle the proof of convergence of the neural network.

\begin{theorem}
    Suppose Assumptions~\ref{as:wp} and \ref{as:3.3} hold, and training reaches \( \mathcal L_{\mathrm{total}}(\theta) \leq \varepsilon \). Then, there exists a constant \( M \) depending only on \(\alpha\) and the Lipschitz constant \(L\) of Assumption~\ref{as:wp} such that
    \begin{equation}
        \sup_{t \in [0,1]} \EE \left [ (N_\theta(t) - y(t))^2 \right ] \leq M\Psi(\hat m, N_q, \varepsilon) 
    \end{equation}
    where \(\Psi(\hat m, N_q, \varepsilon) = C(\varepsilon + \delta_S(\hat m)^2 + N_q^{-2p})\) for a constant \(C\) independent of \(\hat m, N_q, \varepsilon\), with \(\delta_S(\hat m)\) as in Lemma~\ref{lem:3.2}.

    In particular, \(N_\theta \to y\) in \(L^2(\Omega; C[0,1])\) for any sequence such that \(\varepsilon \to 0\), and \(\hat m, N_q \to \infty\).
\end{theorem}
\begin{proof}
    We first quantify the error introduced by the network itself. Define the exact-operator residual \(\tilde r(t)\) comparing the network output \(N_\theta(t)\) against the exact fractional integral operators applied to the learned projections \(b_\theta(t) = \theta_b^\top \mathbf M^\Lambda(t)\) and \(\sigma_\theta(t) = \theta_\sigma^\top \mathbf M^\Lambda(t)\):
    \begin{equation}
        \tilde r(t) = N_\theta(t) - y_0 - \II \left [ b_\theta \right ](t) - \JJ \left [ \sigma_\theta \right ](t).
    \end{equation}
    By Lemma~\ref{lem:3.1}, the deterministic operational matrix is exact: \( \II \left [ b_\theta \right ](t) = \theta^\top_b\mathbb P_\alpha(t)\mathbf M^\Lambda(t) \). By Lemma~\ref{lem:3.2} however, projecting the stochastic integral onto the polynomial basis introduces a truncation error. Let the operational matrix evaluation \(\theta_\sigma^\top \mathbb S_\alpha\mathbf M^\Lambda(t)\) differs from the true It\^o integral \(\JJ \left [ \sigma_\theta \right ](t)\) by a residual process \(\rho_\sigma(t)\) such that \(\EE[\rho_\sigma(t)^2] \leq \Theta^2\delta_S(\hat m)^2\).

    Now, let \(t_i\) be an arbitrary collocation point and define the training residual:
    \begin{equation}
        \hat r (t_i) = N_\theta(t_i) - y_0 - \theta_b^\top \mathbb P_\alpha(t_i)\mathbf M^\Lambda (t_i) - \theta_\sigma^\top \mathbb S_\alpha \mathbf M^\Lambda(t_i).
    \end{equation}
    Substituting this into the exact-operator residual, we have
    \begin{equation}
        \tilde r(t_i) = \hat r(t_i) - \rho_\sigma(t_i).
    \end{equation}
    Squaring and applying the AM--QM inequality, we have
    \begin{equation}
        \tilde r(t_i)^2 \leq 2\hat r(t_i)^2 + 2\rho_\sigma(t_i)^2.
    \end{equation}
    Proceed by taking the expectation, and averaging over all \(N_q\) collocation points:
    \begin{equation}
        \frac 1{N_q} \sum_{i=1}^{N_q} \EE [ \tilde r(t_i)^2 ] \leq \frac 2{N_q} \sum_{i=1}^{N_q} \hat r(t_i)^2 + \frac 2{N_q} \sum_{i=1}^{N_q} \EE [\rho_\sigma(t_i)^2].
    \end{equation}
    The first summation on the right is exactly twice the SDE training loss \(\mathcal L_{\mathrm{SDE}}\) from Equation~\ref{eq:29a}, while the second is bounded as per Lemma~\ref{lem:3.2}. Further note that \(\mathcal L_{\mathrm{SDE}} \leq \mathcal L_{\mathrm{total}} \leq \varepsilon\) by hypothesis. As such, we have
    \begin{equation}\label{eq:50}
        \frac 1{N_q} \sum_{i=1}^{N_q} \EE[\tilde r(t_i)^2] \leq 2\mathcal L_{\mathrm{SDE}} + 2\Theta^2 \delta_S(\hat m)^2 \leq 2\varepsilon + 2\Theta^2 \delta_S(\hat m)^2.
    \end{equation}
    Now, in order to move from a discrete bound at the collocation points to a uniform continuous bound over \([0, 1]\), we apply Lemma~\ref{lem:3.4}:
    \begin{equation}
        \sup_{t \in [0, 1]} \EE [\hat r(t)^2] \leq C_1(\varepsilon + \delta_S(\hat m)^2 + N_q^{-2p}) = \epsilon_1^2
    \end{equation}
    for some constant \(C_1\).

    Next we address the errors that arise from the approximation of the drift and diffusion terms. Define the pointwise deviation functions for each of the projections:
    \begin{subequations}
    \begin{align}
        \xi_b(t) &= b_\theta(t) - b(t, N_\theta(t)) \label{eq:52a} \\
        \xi_\sigma(t) &= \sigma_\theta(t) - \sigma(t, N_\theta(t)) \label{eq:52b}
    \end{align}
    \end{subequations}
    Since the loss functions \(\mathcal L_b\) and \(\mathcal L_\sigma\) are minimized during training, applying Lemma~\ref{lem:3.4} again, their suprema are bounded by constants that depend on training loss and collocation grid resolution:
    \begin{subequations}
    \begin{gather}
        \sup_{t \in [0, 1]} \xi_b(t)^2 \leq \epsilon_2^2, \\
        \sup_{t \in [0, 1]} \EE [\xi_\sigma(t)^2] \leq \epsilon_3^2,
    \end{gather}
    \end{subequations}
    where \(\epsilon_2^2, \epsilon_3^2 = O(\varepsilon + N_q^{-2p})\).

    We now proceed to expand on the global error. Define
    \begin{equation}
        e(t) = N_\theta(t) - y(t).
    \end{equation}
    Subtracting the exact analytic formulation of \(y(t)\) from the definition of \(\tilde r(t)\), we may express the global error \(e(t)\) into five constituent terms as follows:
    \begin{equation}
    \begin{aligned}
        e(t) &= \tilde{r}(t) \\ 
        &- \II \left [ \xi_b\right ](t) \\ 
        &+ \II\left [ b(\cdot, N_\theta(\cdot)) - b(\cdot, y(\cdot)) \right ](t) \\ 
        &- \JJ \left [ \xi_\sigma \right ](t) \\ 
        &+ \JJ \left [ \sigma(\cdot, N_\theta(\cdot)) - \sigma(\cdot, y(\cdot)) \right ](t).
    \end{aligned}
    \end{equation}
    To establish the global bound, we once again square the equation and apply the AM--QM inequality. Take the expectation of both sides and set \(\EE [e(t)^2] = \hat e(t)\).
    \begin{equation}
    \begin{aligned}
        \hat e(t) &\leq 5\EE[\tilde r(t)^2] \\
        &+ 5(\II \left [ \xi_b \right ](t))^2 \\
        &+ 5(\II\left [ b(\cdot, N_\theta(\cdot)) - b(\cdot, y(\cdot)) \right ](t))^2 \\
        &+ 5\EE[(\JJ\left [ \xi_\sigma\right ](t))^2] \\
        &+ 5\EE[(\JJ \left [ \sigma(\cdot, N_\theta(\cdot)) - \sigma(\cdot, y(\cdot)) \right ](t))^2]
    \end{aligned}
    \end{equation}
    We establish bounds with respect to a generic function \(h\). To control the deterministic integrals, apply the Cauchy--Schwarz inequality:
    \begin{subequations}
        \begin{align}
            \left ( \frac 1{\Gamma(\alpha)} \int_0^t (t - \tau)^{\alpha-1}h(\tau)\dd{\tau} \right )^2 &\leq \frac 1{\Gamma(\alpha)^2} \left ( \int_0^t (t - \tau)^{\alpha - 1}\dd{\tau}\right ) \left ( \int_0^t (t - \tau)^{\alpha - 1}h(\tau)^2\dd{\tau}\right ) \\
            &= \frac{t^\alpha}{\alpha\Gamma(\alpha)^2} \int_0^t (t - \tau)^{\alpha-1}h(\tau)^2\dd{\tau}. \label{eq:56b}
        \end{align}
    \end{subequations}
    As for the stochastic integrals, applying It\^o's isometry directly yields
    \begin{equation}\label{eq:57}
        \EE [(\JJ\left [h \right ](t))^2] = \frac 1{\Gamma(\alpha)^2} \int_0^t (t - \tau)^{2\alpha-2} \EE[h(\tau)^2]\dd{\tau}.
    \end{equation}
    Applying Equations~\ref{eq:56b} and \ref{eq:57} in conjunction with Equations~\ref{eq:52a} and \ref{eq:52b} to bound all the integral operators, Equation~\ref{eq:50} to bound the residual, and the Lipschitz bound from Assumption~\ref{as:wp}, we arrive at the following inequality:
    \begin{subequations}
    \begin{align}
        \hat e(t) &\leq 5\epsilon_1^2 + \frac{5t^{2\alpha}}{\alpha^2\Gamma(\alpha)^2} \epsilon_2^2 + \frac{5t^{2\alpha-1}}{\Gamma(\alpha)^2(2\alpha-1)}\epsilon_3^2 \label{eq:58a} \\
        &+ \frac{5L^2}{\alpha\Gamma(\alpha)^2}\int_0^t (t - \tau)^{\alpha-1}\hat e(\tau)\dd{\tau} \label{eq:58b} \\
        &+ \frac{5L^2}{\Gamma(\alpha)^2}\int_0^t (t - \tau)^{2\alpha-2}\hat e(\tau)\dd{\tau}. \label{eq:58c}
    \end{align}
    \end{subequations}
    The monomials in \ref{eq:58a} are monotone and attain maximum value at \( t = 1 \). Furthermore, since \(\alpha \in (1/2, 1)\), the kernel of \ref{eq:58c} is more singular than the kernel of \ref{eq:58b}, allowing us to establish the inequality
    \begin{equation}
        \hat e(t) \leq \Psi + K\int_0^t (t - \tau)^{2\alpha-2}\hat e(\tau)\dd{\tau}
    \end{equation}
    where
    \begin{equation}
        \Psi = 5\epsilon_1^2 + \frac 5{\alpha^2\Gamma(\alpha)^2}\epsilon_2^2 + \frac 5{\Gamma(\alpha)^2(2\alpha-1)}\epsilon_3^2, \quad K = \frac{5L^2}{\Gamma(\alpha)^2} \left ( 1 + \frac 1\alpha \right ).
    \end{equation}

    We now resolve the proof with Lemma~\ref{lem:3.5}:
    \begin{equation}
        \hat e(t) \leq \Psi E_{2\alpha-1}(K\Gamma(2\alpha-1)t^{2\alpha-1}) \leq \Psi E_{2\alpha-1}(K\Gamma(2\alpha-1)).
    \end{equation}
    The second inequality holds due to the monotonicity of the Mittag-Leffler function. Define \(M = E_{2\alpha-1}(K\Gamma(2\alpha-1))\), to arrive at the desired bound:
    \begin{equation}
        \sup_{t \in [0,1]} \EE \left [ (N_\theta(t) - y(t))^2 \right ] \leq M\Psi(\hat m, N_q, \varepsilon).
    \end{equation}
\end{proof}

\begin{remark}
    It is of interest to note that the stochastic projection error depending on \(\hat m\) serves as a bottleneck relative to the deterministic operational matrix. Increasing \(\hat m\) yields diminishing returns once \(\delta_S(\hat m)^2\) drops below \(\varepsilon\) and \(N_q^{-2p}\), at which point resources are better spent reducing the training loss threshold \(\varepsilon\) or increasing the number of collocation points \(N_q\).
\end{remark}

\section{Numerical Experiments}

\lipsum[1]

\subsection{Linear Problems}

\lipsum[2]

\subsection{Nonlinear Problem}

\lipsum[3]

\printbibliography

\begin{appendices}
    \section{Pseudocode} \label{alg:mldnn}
    
    \begin{algorithm}[H]
    \caption{M\"untz-Legendre Neural Network Solver for Caputo SDEs}
    \begin{algorithmic}[1]
    \Require Fractional order $\alpha \in (1/2, 1)$, truncation degree $\hat{m}$, collocation count $N_q$, fine-mesh resolution $n^*$, initial condition $y_0$, drift $b$, diffusion $\sigma$, balancing parameters $\lambda_b, \lambda_\sigma$, learning rate $\eta$, number of epochs $E$
    \Ensure Trained network $N_\theta$ approximating a sample path $y(t, \omega)$ on $[0,1]$
    
    \Statex \textit{Phase~I: Offline Precomputation}
    \State Construct M\"untz exponent sequence $\Lambda \gets \{k\alpha : k = 0, 1, \dots, \hat{m}\}$
    \State Orthogonalize monomials $\{t^{\lambda_k}\}_{k=0}^{\hat{m}}$ to obtain basis $\mathbf{M}^\Lambda(t)$ and coefficient matrix $\mathbf{C}$
    \State Compute diagonal scaling matrix $(\mathbf{D}^\Lambda_\alpha)_{kk} \gets \Gamma(\lambda_k + 1) / \Gamma(\lambda_k + \alpha + 1)$
    \State Assemble deterministic operational matrix $\mathbb{P}_\alpha(t) \gets t^\alpha \mathbf{C} \mathbf{D}^\Lambda_\alpha \mathbf{C}^{-1}$
    \State Sample Brownian path $B(t_j)$ on fine mesh $\{t_j\}_{j=0}^{n^*} \subset [0,1]$ via cumulative Gaussian increments
    \State Compute stochastic integrals on the fine mesh
        \begin{equation} I_i(t) \gets \frac{1}{\Gamma(\alpha)} \int_0^t (t-\tau)^{\alpha-1} M_i^\Lambda(\tau) \dd{B_\tau}, \quad i = 1, \dots, \hat{m} \end{equation}
    \State Compute cross-covariance matrix 
        \begin{equation}\Xi_{ik} \gets \int_0^1 I_i(t) M_k^\Lambda(t)\dd{t}, \quad i, k = 1, \dots, \hat{m} \end{equation}
    \State Compute norm matrix $\boldsymbol \mu \gets \mathrm{diag}\!\left(\norm{M_1^\Lambda}^2, \dots, \norm{M_{\hat{m}}^\Lambda}^2\right)$
    \State Assemble stochastic operational matrix $\mathbb{S}_\alpha \gets \mathbf{\Xi} \boldsymbol \mu^{-1}$
    \State Generate shifted Chebyshev collocation points
        \begin{equation} t_i \gets \frac{1}{2}\left (1 - \cos \frac{(2i-1)\pi}{2N_q}\right ), \quad i = 1, \dots, N_q \end{equation}
    \State Precompute and cache $\mathbf{M}^\Lambda(t_i)$, $\mathbb{P}_\alpha(t_i)\mathbf{M}^\Lambda(t_i)$, and $\mathbb{S}_\alpha\mathbf{M}^\Lambda(t_i)$ at every collocation point
    
    \Statex \textit{Phase~II: Network Training}
    \State Initialize network parameters $\theta_{\text{SDE}}$ and projection vectors $\theta_b, \theta_\sigma \in \mathbb{R}^{\hat{m}}$
    \For{$\mathrm{epoch} = 1$ \textbf{to} $E$}
        \State Evaluate network outputs $N_\theta(t_i)$ at all collocation points
        \State $\displaystyle \mathcal{L}_{\text{SDE}} \gets \frac{1}{N_q} \sum_{i=1}^{N_q} \norm*{ N_\theta(t_i) - y_0 - \theta_b^\top \mathbb{P}_\alpha(t_i) \mathbf{M}^\Lambda(t_i) - \theta_\sigma^\top \mathbb{S}_\alpha \mathbf{M}^\Lambda(t_i) }^2$
        \State $\displaystyle \mathcal{L}_{b} \gets \frac{1}{N_q} \sum_{i=1}^{N_q} \norm*{ \theta_b^\top \mathbf{M}^\Lambda(t_i) - b(t_i, N_\theta(t_i)) }^2$
        \State $\displaystyle \mathcal{L}_{\sigma} \gets \frac{1}{N_q} \sum_{i=1}^{N_q} \norm*{ \theta_\sigma^\top \mathbf{M}^\Lambda(t_i) - \sigma(t_i, N_\theta(t_i)) }^2$
        \State $\mathcal{L}_{\text{total}} \gets \mathcal{L}_{\text{SDE}} + \lambda_b \mathcal{L}_b + \lambda_\sigma \mathcal{L}_\sigma$
        \State Update $(\theta_{\text{SDE}}, \theta_b, \theta_\sigma) \gets (\theta_{\text{SDE}}, \theta_b, \theta_\sigma) - \eta \nabla \mathcal{L}_{\text{total}}$
    \EndFor
    \State \Return $N_\theta$
    \end{algorithmic}
    \end{algorithm}
\end{appendices}

\end{document}